\question 下列关于DRAM和SRAM的说法中，错误的是
Ⅰ．SRAM不是易失性存储器，而DRAM是易失性存储器
Ⅱ．DRAM比SRAM集成度更高，因此读写速度也更快
Ⅲ．主存只能由DRAM构成，而高速缓存只能由SRAM构成
Ⅳ．与SRAM相比，DRAM由于需要刷新，所以功耗较高
\par\twoch{Ⅱ、Ⅲ和Ⅳ}{Ⅰ、Ⅲ和Ⅳ}{Ⅰ、Ⅱ和Ⅲ}{\textcolor{red}{Ⅰ、Ⅱ和Ⅲ和Ⅳ}}
\begin{solution}D。
SRAM和DRAM都属于易失性存储器，掉电就会丢失，Ⅰ错误；SRAM的集成度虽然更低，但速度更快，因此通常用于高速缓存Cache，Ⅱ错误；主存可以用SRAM实现，只是成本高，Ⅲ错误；和SRAM相比，DRAM成本低、功耗低，但需要刷新，Ⅳ错误。
\end{solution}
\question 某SRAM芯片，其容量为512×8位，除电源和接地端外，该芯片引出线的最小数目应该是
\par\twoch{23}{25}{50}{\textcolor{red}{19}}
\begin{solution}D。
容量为512×8位，首先数据线是8位，因为\includegraphics[width=0.64583in,height=0.16667in]{texmath/78dfcc5Cdpi7B3507D25E93D512}，因此地址线为9位，再加上一根读控制线和一根写控制线（可能有些书上的答案还会有电源线、地线等，不要想太多，只算读、写线即可），一共是8+9+2=19，故选D。
\end{solution}
